% 结构版本：D042 五节结构（2.1 测试用例结构分析＋2.2–2.5）。本轮 V011 定稿（2026-09-26 用户对话审定＋三轮 GPT 评审仲裁，AI-12/A048，W-B 窗口）。
% 过程与裁决要点：核料＝三只读子代理（2025A 国一第 2 章写法／题面机制原句逐字摘录／D038 方法侧素材证据状态核验）；
% ①2.5 两段式总路线按 D038 粒度收窄——不出现 REINFORCE/采样解码/精修邻域/"解码代价项"，算法名归 4.4/4.5，场景差异表述升为"场景适配与评价环境"（用户确认；大纲 2.5 素材行字面含"离线 REINFORCE 训练／在线采样解码"，粒度注记待整合窗同步）；
% ②2.4 首句括注"L2 带宽与 DDR 带宽相互独立"保留（GPT 删除建议仲裁驳回、用户改判确认保留：1.2.1 已定稿该事实、1.5.3 有"容量与带宽配置"复述先例、后文"减轻 DDR 压力"需此机制支点；按 D048"信息主要载体＋再引用承担新作用"属合规回指；6.1 仍为 L2 机制唯一展开处）；
% ③"并行度上限"改称"结构并行潜力"（max(W_M,W_V)/CP 无严格上界推导，与 01_context §5.3/§5.4 上限措辞警示一致；T026 后若补推导可改回；大纲 2.1 字样待同步）；
% ④2.2 决策耦合句按场景 A 语境重分配（spill 归 Task 内容与容量、执行顺序管 Task 激活与等待；题面 1.3"同一核内的子图调度顺序影响依赖等待、张量驻留时间和缓存换入换出"句的驻留/换入换出含义主要落场景 B——5.1 起草时对齐口径，勿两章打架）；
% ⑤与 V010（1.3）跨窗对齐：2.1 建图注明求解侧辅助视图（1.3.1"官方为准"原则）；"多消费者张量""分管道计算量"沿用 1.3.2 术语；2.4 重复读取主语由"同一输入/Task"改"同一张量/核心"（1.3.3：L2 复用对象＝同一张量跨核重复 COPY_IN，不限 DDR 图输入）；
% ⑥2.1 数值冻结（D042 数值纪律，% TODO(T026) 源码注释不渲染；T026 出数后从"统计了什么"改写为"数据表明什么"，仅留 2～3 条支撑 4.3 的结论）；TBL-002 未写（仲裁"先无表起草"，定稿按篇幅定去留）；F01 引用句与 figure 环境随图定版补入。
% 素材：题面 1.1/1.2/1.3/1.4/1.5、问题一至三题述、附录 A/C.2/D；无实验数据、无方法名、不挂文献；写法对照 2025A 国一第 2 章（薄分析章：总起定性＋逐问一段＋技术路线图）。
\section{问题分析}\label{sec:analysis}

三个问题共享\textbf{统一的决策骨架}：直接可控的决策同为切图方案与子图调度方案（见 1.2.3 节），方案表示与基本合法性条件一致；差异不在决策对象，而在硬件场景及与之绑定的评价语义——问题二引入核间同步机制，问题三进一步引入全核共享的只读 L2 缓存，评价环境逐问扩展。求解困难的共同来源在于切图、核心分配与执行顺序三层决策的相互耦合：决策空间随操作节点数量与核心数量组合增长，而题设要求在可接受时间内稳定生成高质量方案、避免暴力迭代式搜索。本章 2.1 节先分析测试用例的结构特征，2.2～2.4 节分别讨论三个问题的核心权衡与求解困难，2.5 节在此基础上给出总体建模思路。

\subsection{测试用例结构分析}\label{sec:case-analysis}
% 状态：框架文字落稿（V011），数值冻结待 T026；判定口径随 TBL-012/F22 给出（图或表二选一，选型待真数据，registry F22 状态行待整合窗修正为"工具＋自测就绪、正式数据待 T026"）。
问题分析的对象除机制与约束外，还包括题目给定的测试用例本身：用例的规模、并行潜力与数据共享模式，决定了调度决策空间的基本形态。为刻画与切图决策相关的图结构，本节构建操作级依赖图：收缩张量中转节点，COPY\_IN 与 COPY\_OUT 节点按统计口径处理，建图规则随统计口径统一说明（求解侧辅助视图，正式执行以官方评估程序为准，见 1.3.1 节）。在此图上统计三类信息：一是规模与负载分布，包括操作数、张量数、边数、双管道负载比（由 1.3.2 节分管道计算量核算），并以双管道负载较大者与关键路长度之比刻画结构并行潜力；二是结构类型分布，将用例归入弱连通分量型、汇聚规约型、深链型与其他四类；三是共享数据特征，即多消费者张量（见 1.3.2 节）的数量与规模（对照 L2 容量）。
% TODO(T026, 数值冻结)：官方算例复算并登记 02 后补写——四类结构占比与代表用例、规模与结构并行潜力分布、多消费者张量规模与 L2 容量对照；引出"瓶颈落在可控决策空间"的判断（2.5 与 4.3 的动机链）。外部帖比例与审题报告旧数字（未复算）不入文。

\subsection{问题一分析}
场景 A 无核间同步机制，每个子图独立构成一个 Task，跨子图数据一律经 DDR 中转（见 1.2.3 节）。问题一的困难首先来自三个决策的耦合：切图决定子图粒度、跨子图通信边界以及单个 Task 的缓存工作集，核心分配决定负载分布与哪些数据依赖转为跨核等待，执行顺序则影响 Task 的激活与依赖等待时序；三者共同决定最终的并行收益与搬运代价，题设据此强调不能孤立优化任一环节。核心权衡在于\textbf{并行收益与搬运代价的对抗}：细化切分增加潜在可独立调度的单元，但可能带来更多边界搬运与对同一输入的重复读取；聚合可消除被合并边界上的搬运，却因子图规模增大而加重核内缓存压力，容量不足时评估程序自动插入缓存换入换出。场景 A 中，边界搬运与缓存换入换出同样由 DDR 服务、共同竞争固定的总带宽：提高并行度可能伴随更多并发搬运，从而加剧带宽争用、侵蚀并行收益；粒度变化同时改变两类搬运的规模与并发时序，并进一步影响可并行空间与负载均衡。在满足 1.2.3 节全部合法性条件的前提下，于组合增长的决策空间中高效构造并持续改进方案，构成问题一在求解层面的首要困难。

\subsection{问题二分析}
相较问题一，问题二的方案表示与基本合法性条件沿用不变，场景 B 使核心分配同时承担通信路径选择的作用：相互依赖的子图是否落在同一核心，直接决定其数据通路是核内驻留复用，还是经 DDR 中转并附加同步等待的跨核通路（见 1.2.3 节）。问题二的核心权衡因而是\textbf{数据局部性、负载均衡与缓存压力的三方拉锯}：为利用复用而将子图聚合于同一核心，会压缩跨核并行空间并可能加剧负载不均；为均衡而分散，又把可复用的依赖推向跨核通路；驻留数据持续占用有限的 L1/UB 容量，同核聚合的规模因此受限，题设亦明确要求严格满足容量约束。对求解而言，问题一的方案表示与基本合法性检查可以沿用，但场景 B 的容量要求与同核、跨核代价差异要求搜索策略相应调整——在场景 A 通信结构下形成的偏好，未必仍然成立。

\subsection{问题三分析}
问题三在问题二的基础上引入全核共享的只读 L2 缓存，用于承接多核对共享输入的重复读取（L2 带宽与 DDR 带宽相互独立，见 1.2.1 节）；决策与 Task 组织均沿用问题二。新增的困难在于\textbf{调度与缓存行为的动态耦合}：缓存内容随多核执行过程动态更替，一次读取是否命中，取决于此前哪些数据进入缓存、又是否已被更替，难以在调度决策时静态断定；而切图与核心分配共同决定同一张量被多少个核心重复读取，执行顺序决定这些读取发生的时刻，由此塑造的访问流决定了 L2 所能承接的重复读取比例。L2 的存在同时改变了聚合与分散的权衡：分散在多个核心的相关消费者仍可能经 L2 得到服务、减轻 DDR 压力，但收益取决于命中情况，最终需在官方评估语义下确定；将相关消费者集中到较少核心可减少重复的外部读取，却可能牺牲并行度与负载均衡。因此，问题三的建模须把缓存行为作为调度决策的后果维度加以刻画。

\subsection{总体建模思路}\label{sec:overall-approach}
上述分析确定了统一的建模框架。模型的直接对象是 1.2.3 节的方案表示：切图方案与子图调度方案构成全部可控决策；跨 Task 搬运重建、核内操作排序与共享带宽模拟由官方评估程序承担，总体任务执行时间与总额外数据搬运量等指标均在官方评估语义下定义（见 1.2.3、1.4 节）。由此，本文的建模定位是：在官方评估语义下，为给定计算图与核心数量高效搜索高质量的合法切图与调度方案。从处理流程看，求解统一按输入解析、候选方案构造、合法性检查、评价与优化、结果分析五个环节组织。

从求解方法看，三个问题采用\textbf{统一的两段式路线}：离线段学习引导指派的优先级参数；在线段对每个测试用例依次进行结构感知构造、优先级引导指派与预算内精修，所得方案经合法性检查后交由官方评估程序验证。硬件场景差异不另立总体方法：问题二与问题三的场景语义通过相应的场景适配与评价环境进入同一框架，具体适配方式在第四至六章分别展开。后文进一步构造统一的解析下界，作为候选方案质量的参照（见 4.2.4 节）。
% TODO(F01, 方法冻结/#7 后定版)：图 2-1 全文技术路线图——上层处理流程五环节、下层两段式方法架构＋场景适配落点，只画共用管线与决策边界（D039 分层，算法细节归 F07/F09）；图成后补"全文技术路线如图~\ref{fig:tech-route}所示"句并建 figure 环境。
% TBL-002（三问机制与决策边界对照表）按"先无表起草"落稿（2026-09-26 仲裁）；定稿按篇幅定去留，若留用三列式：问题／相对上一问的新增机制／新增核心权衡（与 TBL-001 字段分工见 05）。落位 2.5。
